% ED26-BM-RAT2021-1-4: exact Goel2021 chapter, printed115--126.
\chapter{Rationality and the Brain}
\label{sec:p19-rationality-1-4}

Goel's chapter asks how a brain can generate inferences while preserving
coherence among beliefs and proposed actions. Its answer is not a single
reasoning organ: it proposes several systems for inference generation,
together with a more general system for handling conflict. The proposed
division depends on conceptual versus logical relations, logical form,
familiar content, presentation modality and determinacy
\cite[pp.~115--124]{Goel2021RationalityBrain}.

The evidence and the interpretation must be distinguished. A contrast in
neural activity is an observation under a particular task and baseline;
assigning it an inferential function is a further claim. Lesions and
stimulation add different constraints, but neither automatically establishes
an isolated, necessary-and-sufficient reasoning module. The chapter's own
note~4 says that the experiments have not clearly separated conflict
\emph{detection} from conflict \emph{resolution}. The account below preserves
that qualification rather than treating two names as two established stages.

The source is the complete 2021 chapter~1.4, not the 2012 Oxford chapter~6.
The two independently inspected empirical anchors are Goel and colleagues'
2007 lesion study and Prado, Chadha and Booth's 2011 imaging meta-analysis.
Other original-study findings are explicitly presented as reports in the
selected chapter. Their missing participant, task, stimulation, analysis or
uncertainty records are not filled with invented details. No new neural
experiment, participant-level reanalysis or population prevalence estimate
is reported here. The reporting interface is
\cref{def:p19-evidence-record}.

\section{Rational action, represented belief and coherence}
\label{sec:p19-r14-criteria}

\begin{hbdefinition}{def:p19-r14-functional-problem}{The source's functional problem}
In Goel's instrumental account, a rational action is a deliberate action
selected from a large or potentially unbounded set of possibilities that advances an organism's
goals coherently with its knowledge and beliefs. A \emph{veridical} belief
fits the world; coherence concerns relations among represented beliefs,
inferences and actions. Neither relationship is identified with a brain
region or with the frequency of a response.

The source distinguishes five requirements: obtaining beliefs that fit the
world; maintaining them in memory; generating inferences; detecting
inconsistency; and updating or revising beliefs. The present neural account
focuses on inference and consistency maintenance. Perception and memory
are required interfaces, not mechanisms fully explained by this account.
\end{hbdefinition}

The thirsty person who believes that a glass of water is on the right of
the desk has a reason to reach right. Reaching left while retaining that
belief is the source's contrasting incoherent action. This example assumes
the thirst-relief goal and the represented location; it is not a universal
claim that leftward movement is irrational. In note~1 the author adds a
philosophical claim, attributed to Cassirer, that antecedent conditions do
not causally suffice for the ensuing rational action. That claim is not a
result of the neural experiments. Note~2 treats epistemic evaluation as
instrumental rationality applied to beliefs and revision. This is Goel's
choice of framework, not a demonstration that all distinctions between
epistemic and practical standards disappear
\cite[pp.~115--116, 124, notes~1--2]{Goel2021RationalityBrain}.

Beliefs, rather than the world directly, mediate the source's action
examples. Believing that a dangerous tiger is under the desk can prompt
flight even when no tiger is there; the same action can be apt if the
belief is true. Conversely, coherent action based on a false location
belief can miss the water. Coherence and veridicality can therefore come
apart. Whether a particular false belief causes harm also depends on the
remaining beliefs, action policy and circumstances; inconsistency alone
does not force a person to enact every conflicting prescription.

\paragraph{Conceptual and logical relations.}
Semantic knowledge supports the connection between being a widow, being a
woman and having had a husband who died, in the source's usage. Logical
connectives and quantifiers constrain patterns involving conjunction,
disjunction, conditionals, universality and existence. Relational expressions
such as greater-than or inside need their own interpretation. The example
that George is taller than Michael and Michael is shorter than George
expresses a relation and its converse. Goel uses such intuitively compelling
connections to explain coherence and regards coherence as primitive. His
description of them as unprovable is not a mathematical impossibility claim:
in a formal language that defines shorter-than as the converse of
taller-than, the connection follows from that definition. The reader must
keep an intuitive psychological starting point separate from the declared
first-order semantics of \cref{sec:fol-semantics-02}
\cite[pp.~115--117]{Goel2021RationalityBrain}.

Induction from observed dogs to unobserved dogs having tails and the causal
explanation of seasons by the Earth's axial tilt illustrate other relations
used in inference. They do not become logical consequences merely because
they are familiar. Sampling, causal assumptions and background knowledge
matter. The source's evolutionary rationale---inference and inconsistency
handling serving truth preservation---is an explanatory proposal, not a
proof that an evolved inductive procedure preserves truth in every case.
The corresponding distinction between evolutionary success and adopted
evaluative criteria is developed in \cref{sec:p19-rationality-1-3}.

\paragraph{The swan revision and its boundary.}
Seeing twelve white swans at Lake Simcoe can lead to the hypothesis that
all swans are white. A later black swan contradicts that universal
hypothesis. The chapter illustrates revision by replacing it with a
most-swans claim. The logically warranted correction from this new
observation is only that \emph{not all} swans are white. It is not that
every swan is nonwhite, and one counterexample does not establish that
most are white. A majority claim needs additional population and sampling
information. This separates successful inconsistency detection from the
adequacy of the particular revision selected
\cite[p.~116]{Goel2021RationalityBrain}.

The chapter deliberately does not explain delusions; note~3 argues for
first understanding ordinary rationality. Its examples consequently cannot
be transferred into a diagnosis of a psychiatric condition. Nor does the
restriction turn neurological patient data into a sample of all ordinary
reasoners \cite[p.~124, note~3]{Goel2021RationalityBrain}.

\section{Inference generation: interpreter, induction and analogy}
\label{sec:p19-r14-generation}

Prefrontal cortex is abbreviated PFC; a Brodmann area is abbreviated BA.
These are anatomical designations, not labels for inference rules. Goel
does not adopt a mental-logic, mental-model, dual-system or probabilistic
theory wholesale before organizing the neural evidence. His proposed
left-PFC interpreter is sensitive to linguistic and conceptual relations;
other frontal and parietal systems contribute to more formal processing,
with a right ventrolateral-PFC contribution to indeterminacy. The proposal
is an organization of the evidence, not a literal implementation of a
specified formal calculus \cite[pp.~117--118]{Goel2021RationalityBrain}.

\paragraph{The split-brain illustration.}
In the Gazzaniga example reported by Goel, a chicken claw is presented in
the right visual field and a snowy scene in the left. The participant's
right hand selects a chicken picture, while the left selects a shovel.
The language-dominant left hemisphere, lacking access to the snow scene,
explains the shovel as useful for cleaning the chicken coop. The proposed
interpreter makes a coherent explanation from the information available
to it, without thereby recovering the actual cause of the other selection.
The visual fields, hand responses, disrupted interhemispheric access and
verbal report are essential to the example. It is not evidence that every
normal explanation is confabulated, nor that the left PFC alone suffices
for rationality. The chapter attributes the demonstration to Gazzaniga
1989 and the interpreter discussion to his 1998 book; it does not supply
the original case sample, recruitment or uncertainty record
\cite[p.~117]{Goel2021RationalityBrain}.

\paragraph{Belief-based induction.}
The description of forty-two-year-old Eve as serious and orderly, enjoying wine and chess and
watching foreign television news can make librarian seem more plausible
than football player. That assessment, attributed to the heuristics
literature, is not an entailment or a verified occupational probability.
Its stereotypic content must not be adopted as a fact about the person.
Similarly, snakes and alligators being cold-blooded does not entail that
all reptiles are cold-blooded. The induction studies of Goel and colleagues
1997 and Goel--Dolan 2004 are reported to engage left dorsolateral PFC,
including BA~9, 8 and 45; the source distinguishes induction from its
deductive comparison by greater left middle-frontal activity around BA~9
\cite[p.~118, fig.~1.4.1a]{Goel2021RationalityBrain}.

The 2019 focal-lesion study varies conclusion believability. Its examples
range from generalizing the residence of one German shepherd to all such
dogs living in D\"usseldorf, through an inference from the glossiness and
moistness of lipstick and fish scales to lipstick containing fish scales,
to the highly believable reptile conclusion. Goel reports less entrenched
beliefs and less acceptance of moderately or highly believable conclusions
following unilateral left BA~9/10 damage. A less believable response is
not automatically a worse inference; acceptance intensity, factual
accuracy and justification are distinct outcomes. The original patient
counts, lesion overlap, response scale and statistical uncertainty remain
unqualified in this source-level report
\cite[p.~118]{Goel2021RationalityBrain}.

\paragraph{Analogical mapping is not one task.}
The chapter juxtaposes three paradigms. Wharton and colleagues' 2000
colored-shape mapping task contrasts analogy with literal identity and
is associated with medial frontal, extensive left prefrontal, anterior
insular and left inferior-parietal activity. Green and colleagues' 2006
word analogy, planet:sun as electron:nucleus, is associated with left
superior frontal BA~9/10 and a frontal-parietal network. Luo and colleagues'
2003 verbal analogy, soldier:army as drummer:band, is reported to involve
bilateral PFC, left temporal cortex and hippocampus. These are different
stimuli and comparisons. Their overlapping activity supports neither
identity of the cognitive operations nor a single analogy locus. No
original-study sample or effect size is supplied by the selected chapter
for these examples \cite[p.~118]{Goel2021RationalityBrain}.

\section{Logical form and belief-laden content}
\label{sec:p19-r14-logical-forms}

Here \(A,B,C\) are temporary schematic names, not measured brain states.
The source contrasts categorical inclusion, transitive comparison,
conditional inference and disjunction. For categorical arguments, all
\(A\)s being \(B\)s and all \(B\)s being \(C\)s entails all \(A\)s
being \(C\)s under the ordinary inclusion interpretation. For a transitive
strict ordering, \(A>B\) and \(B>C\) entails \(A>C\). For the chapter's
deductive conditional benchmark, if Saturday rain rules out Linda's
barbecue and rain occurs, the barbecue is ruled out. For a classical
disjunction, Tom attending with Linda or Mary and not attending with Mary
entails attending with Linda. These benchmarks require their stated
interpretations; an everyday defeasible conditional is not silently
replaced by material implication
\cite[pp.~118, 120]{Goel2021RationalityBrain}.

Goel's broccoli--vegetable--green example has a valid categorical form
even if its universal vegetable premise is false. His London--Paris--Rome
example uses a transitive north-of ordering as a task premise, not a new
geographical measurement. Replacing familiar words by letters preserves
the intended form while changing the content available to a reasoner.
Validity, perceived believability and response accuracy must therefore
be coded separately, as in \cref{sec:p19-rationality-2-3}.

\begin{center}
\begin{tabular}{@{}p{0.09\textwidth}p{0.53\textwidth}p{0.25\textwidth}@{}}
\toprule
Source case & Premises and conclusion & Deductive/content contrast\\
\midrule
H & All broccoli are vegetables; all vegetables are green;
therefore all broccoli are green. & Valid form; source's believable
conclusion.\\
N & No apples are fruit; all fruit contain calories;
therefore no apples contain calories. & Invalid; source's unbelievable
conclusion.\\
O & All apples are fruit; all fruit are poisonous;
therefore all apples are poisonous. & Valid; source's unbelievable
conclusion.\\
P & No apples are fruit; all fruit contain calories;
therefore all apples contain calories. & Invalid; source's believable
conclusion.\\
\bottomrule
\end{tabular}
\end{center}
The table reconstructs the source's congruent cases H/N and incongruent
cases O/P, rather than asserting the truth of its stipulated premises.
For N and P, apples can be disjoint from fruit while being either
calorie-containing or calorie-free. Those two interpretations both satisfy
the premises and separate the proposed conclusions. Believability is a
participant-related content classification; it is not another validity
relation \cite[pp.~120--121]{Goel2021RationalityBrain}.

\paragraph{Form-sensitive neural reports.}
The chapter's quantitative anchor associates categorical reasoning with
left prefrontal BA~9/44, left precentral cortex and basal ganglia;
transitive reasoning with bilateral parietal and dorsal frontal areas;
and its conditional discussion with left dorsal frontal and angular-gyrus
activity. Individual syllogistic studies include Goel 2000, Goel--Dolan
2003 and Reverberi 2012; Reverberi and colleagues' 2009 lesion report
implicates left lateral and superior medial frontal regions in elementary
deduction. The linguistic transitive literature includes Knauff and
colleagues 2003 and Modro\~no and colleagues 2018, and the disjunctive and
conditional examples refer to Eimontaite 2018 and Noveck 2004. These are
source-reported task associations, not evidence that any BA contains a
named logical connective. The source itself regards the smaller conditional
literature as less robust and notes Baggio and colleagues' 2016 left
BA~44/47 findings \cite[pp.~118--120]{Goel2021RationalityBrain}.

\paragraph{Content contrasts and patient constraints.}
For categorical syllogisms, the selected chapter reports left
frontal-temporal involvement for familiar meaningful content and bilateral
PFC, parietal and occipital involvement for corresponding abstract material,
with content controls in Goel 2000 and Goel--Dolan 2003. For familiar
geographical transitive arguments, Goel, Makale and Grafman 2004 report
left parahippocampal and bilateral hippocampal involvement; abstract
counterparts involve bilateral parietal cortex. The controls and baselines
matter: the reported contrast is not a subtraction of one pure reasoning
faculty from another \cite[p.~121, fig.~1.4.1b--f]{Goel2021RationalityBrain}.

Vartanian and colleagues' 2009 frontotemporal-dementia report is a useful
boundary case: relatively spared parietal regions accompany preserved
abstract transitive performance but impaired familiar-content performance.
Goel, Shuren, Sheesley and Grafman's 2004 focal-lesion Wason study reports
similar performance to controls on the arbitrary task but a missing
familiar-content benefit, driven particularly by the left-PFC patients.
Preservation on one version does not imply preservation of every reasoning
capacity. Neither dementia nor a focal-lesion group is randomly assigned
to an anatomical condition. The original samples, lesion distributions,
task scoring and uncertainty for these two reports remain open here
\cite[p.~121]{Goel2021RationalityBrain}.

The rTMS findings reported from Tsujii and colleagues 2010/2011 sharpen
the distinction: left-PFC disruption affects congruent content trials,
not all contentful trials alike. Goel therefore proposes a more restricted
left-PFC contribution through belief bias, conceptual connections and
simple connectives, with other resources needed for complex logical
relations. His possible necessary-and-sufficient interpretation of left
PFC for familiar congruent inference is stronger than the task contrast
alone establishes. A disruption effect can implicate an essential
contribution within a network without showing that the stimulated region
is independently sufficient \cite[p.~121]{Goel2021RationalityBrain}.

\section{The quantitative imaging anchor and its alternatives}
\label{sec:p19-r14-meta}

\paragraph{Corpus, people and comparisons.}
Prado, Chadha and Booth analyzed 28 published peer-reviewed fMRI/PET
studies from 1997--2010, combining 382 participants. Eligibility required
healthy adults, whole-brain activation coordinates in MNI or Talairach
space, and a deductive-reasoning-versus-baseline contrast. Goel's 2007
review provided the initial corpus; PubMed and ScienceDirect searches
covered May 2007--September 2010. Goel--Dolan 2003 was not included because
the required reasoning-versus-baseline coordinates were not reported.
The corpus is not a random sample of all people or all reasoning tasks
\cite[pp.~3484--3488, table~1]{PradoChadhaBooth2011DeductionMeta}.

The selected baselines include sentence comprehension, matching,
nonintegrable premises, fixation, and simpler reasoning. Some contrasts
therefore estimate reasoning against a nonreasoning activity, others the
increment from a simpler to a more complex argument. Material, input
modality, training, imaging method and analyzed phase also vary. The
authors did not have enough studies to examine all those factors
separately. Their three categories are \emph{relational}, \emph{categorical}
and \emph{propositional}. The last includes conditional and disjunctive
inferences, including different forms such as modus ponens and modus
tollens. It must not be renamed a homogeneous conditional-task corpus.

\paragraph{Analysis and reported uncertainty.}
Multilevel kernel density analysis convolves reported peaks with a
10-millimetre-radius sphere and aggregates contrast maps, rather than
counting every peak as an independent participant. Weighting uses the
square root of sample size and an analysis-type multiplier, 1 for
random-effects and 0.75 for fixed-effects analyses. Monte Carlo simulation
with 5,000 iterations supplies the whole-brain threshold; results are
reported at \(p<0.05\), family-wise-error corrected. The target is
spatial consistency of reported activation around a location relative
to the declared spatial null, not the probability that a person uses
a particular cognitive algorithm. No neural effect-size interval or
participant-level reanalysis is manufactured here
\cite[pp.~3485, 3488]{PradoChadhaBooth2011DeductionMeta}.

Across the corpus the reported network is predominantly left
fronto-parietal with basal-ganglia involvement. Category-specific analyses
and direct category contrasts are different results. Relational tasks
consistently involve bilateral posterior parietal and right middle frontal
regions; categorical tasks involve left inferior frontal and basal-ganglia
regions; propositional tasks involve left posterior parietal, precentral
and medial frontal regions. Only left precentral cortex distinguishes
propositional tasks from both other categories in the reported direct
contrasts. For relational tasks, bilateral posterior parietal activity
persists when the five nonlinguistic studies among the eleven relational
studies are removed. Thus visual/nonlinguistic presentation does not
by itself explain the parietal finding
\cite[pp.~3489--3494, tables~2--3]{PradoChadhaBooth2011DeductionMeta}.

\paragraph{What remains a mechanism hypothesis.}
The meta-analysis interprets linear orderings as comparatively easy to
map onto a spatial continuum, while set-inclusion premises allow multiple
configurations and may encourage rule-based representations. Its
propositional result is less decisive: angular-gyrus activity can reflect
verbal as well as spatial processing, medial frontal activity can relate
to abstract-rule maintenance, and precentral activity can reflect
attention or motor demands. The authors explicitly acknowledge that
region-to-process inference lacks a one-to-one mapping and cannot
definitively identify visuospatial versus syntactic computation
\cite[pp.~3492--3494]{PradoChadhaBooth2011DeductionMeta}.

This is a substantive alternative to simply declaring all imaging
findings inconsistent. The source compares meaningful task-dependent
fractionation with methodological differences and with a proposed
left-rostrolateral/medial-frontal deductive core from Monti's studies.
It does not find those proposed core regions consistently across all
tasks; integrating several logical relations in complex tasks is an
alternative explanation for their involvement. Different distributions
of activity constrain a uniform-localization account. They do not,
without a linking theory, prove that no common abstract operation is
implemented in differing networks or that a behavioral model has been
formally refuted. Multiple strategies, individual skills and shared
encoding, memory, control and response operations remain alternatives
\cite[pp.~3494--3495]{PradoChadhaBooth2011DeductionMeta}.

\section{Presentation modality and indeterminacy}
\label{sec:p19-r14-modality}

The chapter contrasts linguistic transitive tasks, usually associated
with parietal involvement, with the pictorial tasks of Fangmeier and
colleagues 2006 and Wendelken--Bunge 2010, associated with anterior or
rostrolateral right BA~10. Waechter and colleagues 2013 propose that
linguistic stimuli must be mapped onto a spatial representation, whereas
the picture already exemplifies the relation. The resulting difference
in representational preparation is an explanatory hypothesis, not a
within-participant modality effect measured by juxtaposing separate
studies. The chapter also reports their lesion comparison: parietal
damage impairs transitive rather than categorical performance. Neither
finding establishes that every pictorial task bypasses parietal cortex
\cite[pp.~120, 122]{Goel2021RationalityBrain}.

The proposed mapping explanation draws on Johnson-Laird's mental models,
Mani--Johnson-Laird's spatial-description work, Cohen and colleagues'
mental-rotation imaging, Goel--Dolan 2001 and Knauff and colleagues 2003.
These interfaces are not interchangeable. In
\cref{sec:p19-mental-rotation}, rotation time can include encoding,
comparison and response components; it does not directly reveal a unique
representation. Likewise, parietal activity during linguistic inference
does not directly reveal a literal internal picture. The current
psychological model account is \cref{sec:p19-rationality-2-3}, not the
model-theoretic meaning of a structure satisfying sentences.

\paragraph{Determinate, inconsistent and unsettled conclusions.}
The source's population example contrasts Toronto larger than Hamilton
and Hamilton larger than Guelph with Toronto larger than each of the
other cities. The first pair determines Toronto larger than Guelph. The
second pair does not determine the Hamilton--Guelph ordering. Under a
strict numerical ordering, take Toronto's value as 3 and the other two
values as 2 and 1 in either order; both assignments satisfy the second
pair but disagree about which of Hamilton and Guelph is larger. Giving
both value 1 supplies a third possibility. This is indeterminacy
\emph{relative to the premises}, not absence of an actual population fact.
It also is not contradiction: both premises can be true
\cite[p.~122]{Goel2021RationalityBrain}.

Rejecting an unentailed conclusion can therefore mean either that it
contradicts the premises or that it remains compatible but unsettled.
The same binary response does not identify the same psychological
operation. Goel reports right-PFC effects for indeterminate relations
from the 2007 patient study, Goel and colleagues 2009 and Brzezicka and
colleagues 2011 imaging, and an indeterminate-syllogism result from
Parsons--Osherson 2001. Ragni and colleagues' 2016 stimulation study
implicates right superior parietal cortex in uncertain spatial reasoning,
but did not test right PFC. It is not a direct replication of the
right-PFC localization \cite[p.~122]{Goel2021RationalityBrain}.

\section{A full primary lesion evidence record}
\label{sec:p19-r14-lesion-record}

\paragraph{People, design and anatomical measurement.}
Goel and colleagues' 2007 study compares nine left-PFC patients, nine
right-PFC patients and 22 age/education-matched controls. Selection
required focal unilateral prefrontal lesions and relatively intact
sensory, motor, language and cognitive functions. Thirteen patients had
Vietnam-era penetrating injuries and were tested in 1999--2002. The five
others had two tumor excisions, an aneurysm and two strokes, with differing
post-injury intervals. Mean ages were 48.50, 54.38 and 49.00 years for
left, right and control groups respectively; mean education was 15.00,
14.67 and 14.84 years. This is a selected clinical comparison, not
randomized damage or a population survey
\cite[pp.~2245--2247, table~1]{GoelEtAl2007CertainUncertain}.

Lesions were traced before behavioral testing from MRI/CT, registered to
a template and compared with approximate BA labels. Damage is not confined
to one BA; BA~10 was involved in eight right and seven left patients,
and the reported group-average losses by BA did not significantly differ.
The printed overlay figure covers fifteen patients, whereas the behavioral
sample is eighteen. The main paper refers to supplementary individual
maps and a supplementary table, which are not supplied by the held
six-page main PDF. The anatomical record cannot be treated as a complete
inspection of all eighteen individual maps.

\paragraph{Task and denominator.}
Participants classified self-paced, visually presented three-term
arguments as valid or invalid using two keys, after a validity explanation
and examples. There were two blocks of 51 trials, with rests after every
25. Half the materials were spatial relations, such as ahead-of; half
were nonspatial relations, such as smarter-than, treated by the study
as orderable relations. The authors collapsed them after reporting
similar accuracy, 71 versus 72 percent, with \(p=0.67\). That decision
does not demonstrate equivalence of their cognitive processes.

The final item counts are 52 valid determinate, 28 inconsistent
determinate and 22 indeterminate, totaling 102. The paper expressly
reports stimulus misclassifications responsible for departures from the
intended balance. Indeterminate trials require ``invalid,'' not a third
uncertainty key. They are not economic risk/reward choices or conditional
probability estimates \cite[p.~2247]{GoelEtAl2007CertainUncertain}.

\paragraph{Results and uncertainty.}
Overall accuracy is reported as 76 percent, SD 15 percentage points.
With education as covariate, the three-group-by-determinacy analysis
reports \(F_{2,33}=6.82\), \(p<0.01\). On determinate items the means
are 72.0 percent for left patients, 78.8 for right patients and 84.7 for
controls; corresponding SDs are 17.6, 11.0 and 14.2. On indeterminate
items the means are 73.5, 56.1 and 79.0, with SDs 19.8, 22.5 and 18.4.
The left-versus-right determinate comparison is explicitly one-tailed
under the stated prior hypothesis. These SDs describe variability of
scores, not confidence intervals for a group difference
\cite[pp.~2247--2248, fig.~2]{GoelEtAl2007CertainUncertain}.

IQ, general and working memory, fluency and naming scores were correlated
with task performance. A second analysis enters six covariates in the
patient-only comparison because the corresponding control scores were
unavailable; it reports an interaction \(F_{1,8}=13.10\), \(p<0.01\).
This adjusted clinical comparison does not establish that every relevant
confound was measured or that all patients used the same strategy.

\paragraph{A separating check on response bias.}
Determinate items mix yes and no responses; indeterminate items all
require no. The authors therefore reanalyze valid, inconsistent and
indeterminate items separately, reporting a lesion-by-item-type
interaction \(F_{2,16}=7.46\), \(p<0.01\). Right versus left accuracy is
56.1 versus 73.5 percent on indeterminate items, 79.8 versus 71.5 on
valid determinate items, but 76.6 versus 75.8 on inconsistent determinate
items. The last comparison is particularly important: it does not fit
an explanation in which right damage simply impairs every no response
or every inconsistency judgment. The source regards response balance
as insufficient to explain the crossover. No new equivalence test or
interval for the near-equal inconsistent means is inferred here
\cite[p.~2248]{GoelEtAl2007CertainUncertain}.

\paragraph{Mechanism and unresolved alternatives.}
The authors propose that right PFC supports maintaining an ambiguous
representation, tempering the left interpreter's premature commitment.
They consider additional working-memory/attention demands as an
alternative, arguing against a simple load account using baseline
scores, covariates, lesion overlap and other imaging observations.
Those arguments narrow alternatives without directly measuring the
postulated ambiguous representation. Their extension to gambling-task
set changes is an explanatory proposal, not a gambling experiment
conducted in these 40 participants. The claimed one-third error
consequence of choosing a specific interpretation also needs an
assumption about how alternatives are selected; it is not a logical
consequence of indeterminacy alone
\cite[pp.~2248--2249]{GoelEtAl2007CertainUncertain}.

The 2021 chapter describes left lesions as producing generalized
impairment and right lesions as selectively impairing incomplete
information. The main 2007 record more specifically supports a
crossover between valid determinate and indeterminate performance,
with similar inconsistent-item means. This preserves the clinical
finding without broadening it into every inferential operation
\cite[p.~122]{Goel2021RationalityBrain}.

\section{Conflict handling is not identified by the final answer}
\label{sec:p19-r14-conflict}

In incongruent syllogisms, an answer cued by conclusion believability
can oppose the validity benchmark. The selected chapter associates
belief-biased responses with ventromedial PFC BA~11/32 and successful
conflict handling with right lateral/dorsolateral PFC, while parietal
resources contribute to formal evaluation. Its cited reports include
Goel 2000, Goel--Dolan 2003, De Neys--Vartanian--Goel 2008,
Prado--Noveck 2007 and Stollstorff--Vartanian--Goel 2012
\cite[pp.~122--123]{Goel2021RationalityBrain}.

A correct answer is compatible with detecting conflict, inhibiting
a belief-driven response and generating a formal response, but it does
not alone observe that entire sequence. A participant can answer by
another strategy, and a wrong answer can follow detection without
successful resolution. This is why the source's explicit detection/
resolution uncertainty matters. Nor does every internally inconsistent
belief set produce conscious conflict on every trial.

\paragraph{Stimulation contrasts.}
As reported by the chapter, right-PFC rTMS impairs incongruent
performance, whereas left-PFC disruption impairs congruent performance
but can improve incongruent performance. The interpretation is that
disrupting a content-sensitive route can reduce irrelevant semantic
interference, permitting other resources to determine the answer.
The intervention contrast does not make impairment a universal decrement
in rationality. For the original Tsujii 2010/2011 reports, sample sizes,
coil placement, stimulation protocol, control/sham conditions, order,
assignment, effect magnitudes and uncertainty are not fully qualified
here; the chapter's reported directions are not a completed intervention
record \cite[p.~123]{Goel2021RationalityBrain}.

\paragraph{Lesion and rule-induction contrasts.}
Caramazza and colleagues' 1976 two-term task asks who is taller and who
is shorter after a taller-than premise. Goel reports widespread left
hemisphere impairment but a right hemisphere deficit for the question
requiring the converse relation. This is a polarity/representation
contrast, not a contradiction in the premise itself. Reverberi and
colleagues' 2005 revised Brixton task separates induction of a rule
from coping with rule conflict; the chapter reports left-PFC impairment
in induction and right-PFC impairment in conflict. These tasks are not
identical to belief--logic incongruence. Their original clinical and
uncertainty records remain unqualified here
\cite[p.~123]{Goel2021RationalityBrain}.

The broader generalization cites Fink 1999, Picton 2006, two Vallesi
2007 reports, and Marinsek 2014. They contribute different intention,
timing, temporal-preparation and hemispheric paradigms, not repeated
instances of the same deductive task. Hamburger and colleagues 2018
and Knauff's spatial-reasoning account supply a related case in which
visual imagery can interfere with spatial reasoning; they do not
identify every visual or right-PFC operation as conflict detection.
Without their task-specific records, these citations support attribution
of the source's generalization rather than a new cross-paradigm effect
estimate \cite[p.~123]{Goel2021RationalityBrain}.

\section{Localization, causal interpretation and neural plausibility}
\label{sec:p19-r14-limits}

The chapter's figure collects nine panels from distinct studies. Panels
a--h concern induction, meaningful syllogisms, abstract syllogisms,
familiar and abstract transitive relations, pictorial integration and
indeterminate relations; panel~i concerns conflict handling. It is not
a factorial experiment in a single sample. Its caption labels the
conflict panel BA~44/45, while the later text gives BA~45/46. The
qualitative right lateral/dorsolateral attribution can be preserved,
but those differing anatomical details cannot be silently reconciled
into a new precise locus without the original imaging record
\cite[p.~119, fig.~1.4.1; p.~123]{Goel2021RationalityBrain}.

Three inferential boundaries are consequential. First, a measured
signal can covary with a condition containing encoding, representation,
memory, inference, control and motor response; it is not a direct readout
of one process. Second, disrupting a region can affect connected
systems and a particular task contribution; an effect does not show
that the region alone computes the inference. Third, lesion groups
can differ in history, extent, preserved resources and compensation;
adjustment and dissociations constrain those alternatives but do not
turn the design into random anatomical assignment. These are reasons
to triangulate task contrasts, behavior, imaging and intervention,
not reasons to discard neural evidence.

The mapped 2012 Oxford chapter, Morrison--Knowlton's \emph{Neurocognitive
Methods in Higher Cognition}, concerns what neural measurements and
interventions can establish about a reasoning mechanism. Its mapped
questions include behavioral underdetermination, neuropsychological
cases, EEG/fMRI measurements, TMS and prefrontal cognition. That is a
methodological comparison interface, not another edition of Goel's
2021 multi-system account. Its full chapter is not qualified by the
present source corpus, so no uninspected Oxford argument or completed
coverage credit is supplied here. The present lesion and imaging
records illustrate the questions at that interface without replacing
the missing Oxford treatment.

The existing P19 dual-process account requires specification of
representations, computations, engagement and interactions
(see \cref{def:p19-r25-process-account,def:p19-r25-interactions}). Goel's
interpreter/formal/conflict organization is not a direct anatomical
identification of Type~1 and Type~2. A belief-based route can be correct,
a resource-demanding route can err, and multiple neural resources can
contribute to one response. Likewise, the task/construct/measure
distinction in \cref{def:p19-empirical-question} rules out declaring a
participant rational from an activation label. An empirical neural
constraint does not establish the normative benchmark against which
the response is judged.

\section{Using the account for an actual reasoning question}
\label{sec:p19-r14-use}

To evaluate a proposed neural explanation, first state the argument
interpretation and response criterion; then separate form, content,
modality, determinacy and conflict. Record the people, comparison,
measurement or perturbation, baseline and reported uncertainty. Ask
which competing mechanism would produce the same score or signal.
The success case is the 2007 lesion contrast: the indeterminate deficit
is compared with valid and inconsistent determinate items, so a simple
no-response explanation can be tested against the actual pattern.
The failure case is to infer a unique uncertainty mechanism merely
from a right-frontal activation, or to infer a general loss of
rationality merely from a focal clinical deficit.

Goel's conclusion is a provisional multi-system account: linguistic/
conceptual inference, form- and modality-sensitive visuospatial
resources, indeterminacy tolerance and more general conflict handling
cooperate in maintaining a useful belief network. The source expects
the anatomical and functional details to change with further evidence.
Its broad rejection of a unitary reasoning module is an empirical
interpretation of the reviewed studies, not an impossibility theorem
about all common computations. Correctness of a formal inference,
accuracy of a belief, observed clinical performance and rationality
under adopted goals remain four different questions
\cite[pp.~123--124]{Goel2021RationalityBrain}.
